% problem-000198 1 镓氯化合物与团簇
\section*{1. 镓氯化合物与团簇}
\textit{下列为分问。}

\subsection*{1-1}
半导体⼯业中通过刻蚀制造微纳结构。GaN是重要的半导体材料，通常采⽤含氯⽓体在放电条件下进⾏刻蚀。写出利⽤ $\mathrm { A r } { - } \mathrm { C l } _ { 2 }$ 混合⽓体放电刻蚀GaN的化学⽅程式。

\subsection*{1-2}
⾦属镓熔点很低但沸点很⾼，其中存在⼆聚体 $\mathrm { G a _ { 2 } } _ { \mathrm { { c } } }$ 。1990年，科学家将液态Ga和 $\mathrm { I } _ { 2 }$ 在甲苯中超声处理，得到了组成为GaI的物质。该物质中含有多种不同氧化态的Ga，具有两种可能的结构，分⼦式分别为$\mathrm { G a } _ { 4 } \mathrm { I } _ { 4 } \left( \mathrm { A } \right)$ 和 $\mathrm { G a } _ { 6 } \mathrm { I } _ { 6 } \left( \mathrm { B } \right)$ ，⼆者对应的阴离⼦分别为C和D，两种阴离⼦均由Ga和I构成且其中所有原⼦的价层均满⾜8电⼦。写出⽰出A和B组成特点的结构简式并标出Ga的氧化态，画出C和D的结构。

\subsection*{1-3}
GaI常⽤于合成低价Ga的化合物。将GaI与 $\mathrm { A r ^ { \prime } L i }$ （Ar′基如图所⽰，解答中直接采⽤简写 $\mathrm { A r ^ { \prime } } )$ ）在−$7 8 ~ ^ { \circ } \mathrm { C }$ 的甲苯溶液中反应，得到晶体E，E中含有2个Ga原⼦；E在⼄醚溶液中与⾦属钠反应得到晶体F，X射线晶体学表明，F中的Ga-Ga键⻓⽐E中短0.028nm。关于F中Ga-Ga的键级历史上曾有过争议，其中⼀种观点认为，F中的Ga价层满⾜8电⼦。基于该观点，画出E和F的结构式。

（图：）
